%!TEX program = xelatex
%!TEX root = ../all.tex

\chapter{Decision Trees}

A \alert{decision tree} is a supervised learning model designed for either classification or
regression tasks. At its core, the idea is disarmingly simple: rather than fitting a single
global function to the data, a decision tree builds a sequence of \emph{if--then} decisions
that progressively refine our knowledge about the output. Starting from the full dataset, the
tree repeatedly partitions the instance space into smaller and smaller regions, each one more
homogeneous with respect to the target variable than the one it came from. In this sense, a
decision tree can be understood as a \emph{recursive partition} of the feature space, where
each internal decision reduces ambiguity and each terminal region provides a final prediction.

What makes decision trees particularly appealing as a modeling tool is the combination of
predictive power and interpretability. Each internal node in the tree corresponds to an
explicit, human-readable test on one of the input features, and a prediction for a new
instance can be explained simply by tracing the sequence of tests that the instance satisfies
from the root down to a leaf. This transparency is often a decisive practical advantage,
especially in domains such as medicine, law, or finance, where the rationale behind a
model's prediction carries as much weight as the prediction itself. Unlike many other
machine learning models, a decision tree does not hide its logic inside a large collection
of numerical parameters: its reasoning is laid out openly as a tree of conditions that any
practitioner can inspect, question, and communicate to a non-technical audience.

\section*{Structure and components}

Formally, a decision tree is a rooted hierarchical structure composed of nodes and directed
edges. The \textbf{root node} sits at the top of the hierarchy and represents the entire
training set before any partition has taken place: it is the starting point from which all
decisions flow. From the root, the tree expands through a sequence of \textbf{internal
nodes}, each of which implements a decision rule based on the input features. The rule at an
internal node tests some property of the instance currently being classified---for example,
whether a particular feature exceeds a threshold---and routes the instance to one of the
node's children depending on the outcome of the test. This routing continues recursively
until the instance reaches a \textbf{leaf node} (also called a terminal node), which
represents a region of the feature space where no further splitting is performed and where a
final prediction is issued.

An important practical convention is that each internal node is typically associated with a
partition induced by a \emph{single} attribute. This univariate splitting constraint is not
strictly necessary from a mathematical standpoint---one could in principle test arbitrary
functions of several features at once---but it plays a crucial role in making the resulting
model interpretable. When a split depends on exactly one feature (for instance, ``petal
length $\leq 2.45$''), the meaning of each branch is immediately understandable to a human
reader without any further explanation. As the tree grows deeper, each path from the root
to a leaf accumulates a conjunction of such simple conditions, and the corresponding leaf
therefore describes a clearly delineated subregion of the feature space: the set of all
instances satisfying every condition along that path simultaneously.

The nature of the prediction stored at each leaf depends on the task. In
\textbf{classification}, a leaf may store a single \emph{class label}, yielding a
deterministic prediction for every instance that falls into the corresponding region.
Alternatively, it may store the \emph{empirical distribution of class labels} observed in
that region during training, allowing the model to output calibrated probability estimates
rather than hard assignments---a useful capability when downstream decisions depend on
prediction confidence. In \textbf{regression}, the leaf stores a numerical summary of the
target values observed in that region, most commonly their mean, which is then returned as
the predicted output for any new instance landing in the leaf.

It is worth pausing to appreciate how naturally the tree structure encodes the idea of
divide-and-conquer problem solving. The root node faces the hardest problem: predicting the
output for the entire, presumably heterogeneous, dataset. By choosing an informative
splitting criterion, it divides this hard problem into two (or more) easier sub-problems,
each involving a more homogeneous subset of the data. Each child node then applies the same
logic recursively to its own subset, further reducing heterogeneity at every level. When the
subsets become sufficiently pure---either because they contain instances of a single class,
or because a stopping criterion is met---the recursion terminates and leaf nodes are created.
The depth of the tree and the number of leaves thus reflect the complexity of the learned
function: a shallow tree imposes a coarser partition and a simpler decision boundary, while
a deep tree can capture finer distinctions at the risk of overfitting to the training data.

\subsection*{A visual example}

The figures below illustrate how a decision tree is grown, step by step, on the classic Iris
dataset, which contains measurements of petal and sepal dimensions for flowers of three
species. The visualizations serve a dual purpose: they show both the tree structure itself
and the corresponding geometric partition of the feature space induced by the splits. Reading
these two representations in parallel is essential for building intuition about what a
decision tree actually computes.

\begin{figure}[ht]
\begin{center}
\includegraphics[width=.7\linewidth]{iris.png}
\caption{The Iris dataset}\label{f:iris}
\end{center}
\end{figure}

The scatter plot in Figure~\ref{f:iris} shows the raw data. Even without any model, one can see that the
three species form clusters in the feature space, though some overlap exists, suggesting
that a small number of well-chosen splits might already achieve good separation.

\begin{figure}[ht]
\begin{center}
\includegraphics[width=.4\linewidth]{dec_tree_0.png}
\centerline{\textit{Initial decision tree with first split}}
\bigskip
\includegraphics[width=.7\linewidth]{dec_tree_regions_0.png}
\centerline{\textit{Corresponding partition of the feature space}}
\caption{After the first split}\label{f:split1}
\end{center}
\end{figure}

The first split introduces a single axis-aligned cut in the feature space, dividing the
instances into two groups (Figure \ref{f:split1}). One of the resulting regions may already be pure enough to
constitute a leaf, while the other still contains a mixture of classes and must be split
further. The corresponding tree has a root node and two children: one leaf and one internal
node awaiting refinement. Notice how the partition of the feature space at this stage
consists of just two half-planes separated by a single decision boundary.

\begin{figure}[ht]
\begin{center}
\includegraphics[width=.4\linewidth]{dec_tree_1.png}
\centerline{\textit{Decision tree after second split}}
\bigskip
\includegraphics[width=.7\linewidth]{dec_tree_regions_1.png}
\centerline{\textit{Updated feature space partition}}
\caption{After the second split}\label{f:split2}
\end{center}
\end{figure}

The second split refines the partition by introducing a further cut within the impure
region identified at the previous step (Figure~\ref{f:split2}). The feature space is now divided into three
regions, and the tree has grown one level deeper in the branch that required further
discrimination. At each step, the algorithm selects the split that maximally reduces the
impurity (or heterogeneity) of the resulting subsets, according to some criterion such as
the Gini index or information gain.

\begin{figure}[ht]
\begin{center}
\includegraphics[width=.6\linewidth]{dec_tree.png}
\centerline{\textit{Final decision tree}}
\bigskip
\includegraphics[width=.7\linewidth]{dec_tree_regions.png}
\centerline{\textit{Final partition of the feature space}}
\caption{After the third split}\label{f:split3}
\end{center}
\end{figure}

The final tree and its associated partition illustrate the end result of the growing
process (Figure~\ref{f:split3}). The feature space has been divided into a collection of rectangular regions (a
direct consequence of using univariate, axis-aligned splits at each node), and each region
is assigned the class label that was most common among the training instances it contains.
The boundaries between regions are rectilinear, reflecting the fact that every split is
determined by a threshold on a single feature. This axis-aligned structure is one of the
most recognizable geometric signatures of decision tree models, and it distinguishes them
from classifiers such as linear discriminants or support vector machines, whose decision
boundaries are oblique hyperplanes.

Comparing the tree diagram with the partition plot reinforces a key point: every path from
the root to a leaf in the tree corresponds to exactly one rectangular cell in the partition,
and conversely, every cell in the partition corresponds to exactly one leaf in the tree. The
two representations are therefore dual to each other, and understanding this duality is
essential for reasoning about how a decision tree generalizes to new data, where it succeeds,
and where it can fail.

\subsection*{Decision Tree: classification process}

Once a tree has been learned from training data, classifying a new instance amounts to a
deterministic \emph{top-down traversal} from the root to a leaf. Let
$\z=(z_1,\ldots,z_d)^T$ denote an input instance with $d$ features. The prediction process
begins at the root node, which represents the full feature space before any decision has
been made. At each internal node encountered along the way, the model evaluates the node's
decision rule and selects the outgoing branch that corresponds to the outcome of that rule,
thereby restricting attention to a progressively smaller subregion of the feature space. The
instance carries no memory of the path it has already traveled: at each step, only the
current node's rule matters.

More explicitly, suppose an internal node is associated with a feature $x_i$ and a decision
function $f_i$. The evaluation of the node consists of computing $f_i(z_i)$ and using its
value to select the next child node. A very common case is a threshold test applied to a
continuous feature, where the node stores a threshold $\theta_i$ and implements the binary
decision:
\begin{myequation}
\text{if } z_i \leq \theta_i \text{ go left,\quad otherwise (if } z_i > \theta_i \text{) go right.}
\end{myequation}
Each such test can be interpreted geometrically as a cut of the current region along an
axis-aligned hyperplane, and statistically as a way of isolating subpopulations whose
members tend to share similar labels. The sequence of cuts accumulated from the root down to
any given node defines the rectangular cell of the feature space to which that node
corresponds.

The traversal continues, node by node, until a leaf is reached. At that point the tree
returns either a class label, representing a deterministic prediction, or a probability
distribution over classes, from which a final label can be derived if needed---for instance
by selecting the most probable class. Both output modes will be discussed more carefully when
we consider how leaf nodes are assigned their predictions.

From a computational standpoint, because the traversal inspects exactly one node per level
of the tree, the time required to produce a prediction is proportional to the height $h$ of
the tree, that is $O(h)$. This makes inference extremely efficient: even when the training
set is very large and the learned tree is complex, classifying a new instance requires only a
small number of simple comparisons. This property is one of the practical strengths of
decision trees relative to models such as nearest-neighbor classifiers, whose prediction cost
grows with the size of the training data.

\subsection*{Decision Tree: construction algorithm}

Understanding how a decision tree is built helps clarify why it is both powerful and
potentially prone to overfitting. Tree construction proceeds recursively, starting at the
root and repeatedly selecting a split that partitions the current set of training instances
into more homogeneous subsets. The most common approach is \emph{greedy and top-down}: at
each node one chooses the locally best split according to some criterion, applies it
immediately without reconsidering earlier choices, and then continues independently on each
resulting child subset. The greedy nature of this strategy means that the algorithm does not
search for the globally optimal tree---a problem that is computationally intractable in
general---but instead builds a reasonable tree efficiently by making a sequence of locally
good decisions.

Operationally, one starts with the entire training set at the root. For a generic node
associated with a subset $S$ of the training data, the algorithm first decides whether $S$
should become a leaf or whether it should be split further. If splitting is chosen, the
algorithm selects the feature and split rule that best improves the purity of the resulting
subsets, partitions $S$ into child subsets accordingly, and then applies the same procedure
recursively to each child. When a node is declared terminal, it is assigned prediction
information: a class label, class distribution, or numerical value, depending on the task at
hand.

At every step of this process, two conceptual questions are always present, even if a
library implementation hides them behind a single function call. The first is \emph{how to
partition}: which feature should be tested, and with which threshold or grouping rule, in
order to obtain the most useful separation of classes or the greatest reduction of regression
error? The second is \emph{when to stop}: at what point should the current region be
accepted as sufficiently informative and committed to a prediction, rather than split further
at the risk of fitting noise rather than signal? These two questions are addressed in turn
in the sections that follow.

\subsection*{Decision Tree: selecting the best split}

The selection of the best split at each node is guided by the notion of \alert{impurity}.
Intuitively, a node is ``good'' when the labels of the instances that fall into it are mostly
the same---high homogeneity---and ``bad'' when it contains a mixture of different classes---
high heterogeneity. Impurity measures formalize this intuition: they assign a real number to
a class distribution, small for concentrated (pure) distributions and large for uniform
(maximally mixed) ones. The goal of the splitting step is therefore to find the partition
that reduces impurity the most, moving from a more heterogeneous parent node to more
homogeneous children.

Before discussing specific impurity measures, it is useful to state precisely what
properties a good impurity function should satisfy, and why each property is natural.

\subsubsection*{Impurity measure: formal definition}

Consider a node containing instances drawn from $k$ possible classes, and let
$p = (p_1, \ldots, p_k)$ denote the empirical class distribution at that node, where $p_i$
represents the fraction of instances belonging to class $i$. An impurity measure is a
function $\phi: p \mapsto \mathbb{R}$ intended to quantify how mixed the distribution is.
Several natural requirements characterize what we expect from such a function.

To begin with, impurity should never be negative:
\begin{myequation}
\phi(p) \geq 0 \quad \text{for all } p.
\end{myequation}
This is a basic sanity condition: impurity is a quantity to be minimized, and having a
non-negative baseline makes the notion of ``zero impurity'' (perfect purity) well-defined.
Second, impurity should be minimal when the distribution is perfectly concentrated on a
single class, that is, when there exists an index $i$ such that $p_i = 1$. For most common
choices, $\phi(p) = 0$ in this case, reflecting the fact that a node containing instances
of only one class requires no further splitting. Third, impurity should be maximal when all
classes are equally represented, that is when $p_i = 1/k$ for all $i$, since this is the
situation of greatest uncertainty about the label of a randomly drawn instance.

Two additional properties are typically required. Symmetry states that impurity should not
depend on the ordering or naming of the classes:
\begin{myequation}
\phi(p) = \phi(p')
\quad
\text{for every } p' \text{ obtained by permuting the components of } p.
\end{myequation}
This is natural because the classes themselves carry no intrinsic ordering: relabeling them
should not change how mixed the distribution is. Smoothness, meaning differentiability of
$\phi$ on the interior of the probability simplex, is also a convenient property, especially
when one wants to study the behavior of impurity-based criteria analytically or to optimize
impurity-related objectives with gradient methods:
\begin{myequation}
\phi(p)\ \text{is differentiable on the interior of the simplex.}
\end{myequation}

Taken together, these conditions encode a simple and compelling interpretation: impurity
measures how hard it is to correctly guess the class label of a randomly drawn instance from
the node. If almost all instances share the same label, guessing is easy and impurity is
close to zero. If classes are balanced, no label is a better guess than any other, and
impurity is large. The impurity function thus quantifies, in a single number, how much
uncertainty about the label remains at a given node.

\subsubsection*{Goodness of split}

Given a labeled set $S$ of instances, let $p_S$ denote the class-probability vector
associated with $S$, computed as empirical frequencies. Suppose we apply a split function
$f: S \mapsto \{1, \ldots, r\}$ that partitions $S$ into $r$ disjoint subsets:
\begin{myequation}
s_i = \{x \in S \mid f(x) = i\},
\end{myequation}
so that $S = s_1 \cup \cdots \cup s_r$.

The central idea is to compare the impurity of the parent node with the \emph{expected
impurity} after the split. The expected impurity after splitting is computed as a weighted
average of the impurities of the child nodes, with weights proportional to the fraction of
data that falls into each child. This leads to the following definition of
\textbf{goodness of split}:
\begin{myequation}
\Delta_\phi(S, f) = \phi(p_S) - \sum_{i=1}^r \frac{|s_i|}{|S|} \phi(p_{s_i}).
\end{myequation}
The quantity $\Delta_\phi(S, f)$ measures the reduction in impurity achieved by splitting
$S$ according to $f$. A larger value of $\Delta_\phi$ means that the children are, on
average, substantially purer than the parent, so the split is considered better. When
searching for the best split at a node, the algorithm evaluates $\Delta_\phi(S, f)$ for all
candidate split functions $f$ and selects the one achieving the highest value.

The weighted average $\sum_{i=1}^r \frac{|s_i|}{|S|} \phi(p_{s_i})$ in the definition
deserves a moment of attention. The weighting by $|s_i|/|S|$ ensures that large child nodes
contribute more to the expected post-split impurity than small ones. This is desirable: a
split that achieves excellent purity in a tiny subset while leaving the majority of instances
in a mixed node should not score as well as a split that achieves moderate purity improvement
across all children. The weighted average correctly captures this asymmetry.

\subsubsection*{Practical split functions}

While the definition above allows for very general split functions, practical decision trees
typically restrict attention to splits based on a \emph{single feature}. This restriction is
closely tied to interpretability and computational feasibility: it limits the search space of
candidate splits and ensures that each decision rule can be stated as a simple condition on
one variable, making the tree easy to explain.

For a \textbf{discrete} (categorical) feature, one can create a multiway split by assigning
one branch to each distinct value of the feature, or alternatively group the values into two
sets to obtain a binary split. In the multiway case, the split function simply maps each
categorical value to its own child index. In the binary case, the algorithm must find the
best partition of the set of categories into two groups, which for small cardinalities can
be done by exhaustive enumeration. The binary version is sometimes preferred when the feature
has high cardinality, since a multiway split would produce many small, statistically
unreliable children.

For a \textbf{continuous} feature, splits are defined by one or more thresholds. The most
common and computationally tractable choice is a single threshold $\theta$, producing the
binary partition:
\begin{myequation}
x_i \leq \theta
\quad \text{vs} \quad
x_i > \theta.
\end{myequation}
Candidate thresholds are typically chosen from the set of observed feature values. A standard
and exact procedure is to sort the training instances by the feature value, consider all
midpoints between consecutive distinct values as candidate thresholds, evaluate $\Delta_\phi$
at each candidate, and select the threshold achieving the highest goodness of split. The
total cost of this search for a single feature is $O(n \log n)$ due to sorting, where $n$ is
the number of instances at the current node.

\subsubsection*{Entropy and information gain}

A particularly important and theoretically well-motivated impurity measure comes from
information theory: \alert{entropy}. Entropy quantifies the uncertainty present in a
probability distribution. When applied to the class distribution at a node, it captures how
unpredictable the class label of an instance reaching that node is: a node with high entropy
is one where many different classes are represented in comparable proportions, while a node
with low entropy is one where most instances share the same label.

For any set $S$ of instances drawn from $k$ classes, with $S_i$ denoting the subset
belonging to class $i$, entropy is defined as
\begin{myequation}
H_S = -\sum_{i=1}^k \frac{|S_i|}{|S|} \log_2 \frac{|S_i|}{|S|}.
\end{myequation}
The base-2 logarithm is a conventional choice that gives entropy the unit of bits, matching
the classical information-theoretic interpretation: $H_S$ equals the expected number of
binary questions one would need to ask (under an optimal strategy) in order to identify the
class of a randomly chosen instance from $S$.

The extreme values of entropy confirm that it behaves as desired. When the node is perfectly
pure and all instances belong to the same class $i$, then $|S_i|/|S| = 1$ and
$|S_j|/|S| = 0$ for $j \neq i$. Using the convention $0 \log_2 0 = 0$, we obtain $H_S = 0$:
there is no uncertainty and no information is needed to determine the class. Conversely, when
all $k$ classes are equally represented, $|S_i|/|S| = 1/k$ for every $i$, and
$H_S = \log_2 k$, which is the maximum possible value. This upper bound grows with the
number of classes, reflecting the fact that more classes means more potential uncertainty.

When entropy is used as the impurity function $\phi$, the goodness-of-split criterion
$\Delta_\phi(S, f)$ specializes to what is known as \alert{information gain}. For a split
function $f$ that partitions $S$ into subsets $s_1, \ldots, s_r$, information gain is:
\begin{myequation}
IG(S, f) = H_S - \sum_{j=1}^r \frac{|s_j|}{|S|} H_{s_j}.
\end{myequation}
The interpretation

\subsection*{Pruning: overview and motivation}

The discussion above highlights a fundamental tension in decision tree learning: trees can
become very expressive by growing deep and creating fine-grained partitions, but this
expressiveness comes with a high risk of \emph{overfitting}. A decision tree that is allowed
to grow without restraint will eventually create one leaf per training instance, achieving
zero training error while learning nothing transferable to new data. Unlike models whose
complexity is tuned through continuous parameters---where regularization can be applied
smoothly---a decision tree's complexity emerges from discrete splitting decisions, making the
bias--variance trade-off both essential to manage and genuinely difficult to control.

\paragraph{The core problem}

The key difficulty is that it is rarely obvious, during the growth phase, when the tree
should stop adding splits. Two extreme strategies illustrate the dilemma clearly.

On one end, enforcing strict stopping criteria---known as an \emph{early stopping}
strategy---produces small, shallow trees. Such trees are easy to interpret and often
statistically stable across different training samples, but they may suffer from high bias
and underfit the data: important patterns can require several levels of splitting before
becoming apparent, and stopping too early may prevent the tree from discovering deeper
regularities. This is sometimes described as a ``horizon effect'': a split that looks only
moderately useful when evaluated locally might unlock a sequence of highly informative splits
further down the tree, which early stopping would never allow the algorithm to reach.

On the other end, adopting loose stopping criteria---a \emph{late stopping} strategy---lets
the tree continue splitting until it has captured almost every irregularity of the training
set. Training error can become extremely small, sometimes even zero, but many of the learned
distinctions reflect noise in the specific training sample rather than genuine structure in
the underlying distribution. The resulting model has high variance: it is sensitive to small
fluctuations in the training data and generalizes poorly to unseen instances. Furthermore,
deeply grown trees are often large enough to be practically uninterpretable, undermining
one of the main advantages of the decision tree family.

\paragraph{The pruning solution}

Pruning addresses this dilemma with a pragmatic ``grow first, simplify later'' philosophy.
Rather than trying to guess the correct stopping time during growth---a problem complicated
by the horizon effect---one intentionally grows a large tree that has the opportunity to
capture all potentially useful regularities, and then systematically removes branches that
do not contribute to generalization. The strategy decomposes naturally into two phases: a
\textbf{growing phase}, in which the tree is expanded using loose stopping criteria,
favoring coverage and flexibility, and a \textbf{pruning phase}, in which subtrees that
appear to model noise rather than robust structure are deleted, retaining only the parts that
improve expected performance on unseen data.

The pruning phase can estimate generalization in several ways: by holding out a dedicated
validation set and measuring accuracy directly, by applying a statistical correction to the
training error to account for optimism, or by penalizing model complexity through a
regularization term. Different pruning methods reflect different choices along these
dimensions. The following sections introduce and discuss the most important approaches, each
of which has influenced the design of major decision tree implementations.

\subsubsection*{Reduced Error Pruning (REP)}

Reduced Error Pruning is perhaps the conceptually simplest pruning
method and provides a useful baseline against which more sophisticated approaches can be
compared. The algorithm requires the data to be partitioned into three disjoint sets: a
training set used to grow the initial tree, a \emph{pruning set} (also called a validation
set) used to evaluate pruning decisions, and an optional test set for final evaluation. The
pruning set plays a central role: all decisions about which subtrees to remove are made by
directly measuring their effect on the accuracy of the tree on this held-out data.

The algorithm proceeds bottom-up through the internal nodes of the fully grown tree. At each
internal node, the algorithm asks: what would happen if this node were replaced by a leaf
that predicts the majority class among the training instances that reached it? It then
evaluates both the original subtree and the candidate leaf on the pruning set, and replaces
the subtree with the leaf if doing so does not decrease accuracy on the pruning set. The
process continues until no further pruning improves or maintains accuracy.

Several properties of REP deserve attention. First, it is exact in the sense that it
directly optimizes the criterion of interest (accuracy on held-out data) rather than
approximating it. Second, it is computationally efficient: a single bottom-up pass through
the tree suffices, and each pruning decision requires only a simple comparison of
error counts on the pruning set. Third, and most importantly, REP tends to produce small
trees that generalize well when the pruning set is sufficiently large, because it directly
measures generalization rather than estimating it indirectly.

The main limitation of REP is its dependence on the size of the pruning set. If the
available data are scarce, allocating a substantial fraction to pruning leaves too little
for growing the initial tree, and the resulting pruned tree may be overly aggressive in
removing branches that would have been useful with more training data. Conversely, a small
pruning set produces noisy estimates of validation accuracy, leading to unreliable pruning
decisions. For this reason, REP is most attractive when large amounts of labeled data are
available, and it is typically used as a reference method rather than as the preferred
approach in data-scarce settings.

\subsubsection*{Cost Complexity Pruning (CCP) and the CART framework}

Cost Complexity Pruning, sometimes called weakest-link pruning, is the pruning strategy
employed in the CART (Classification and Regression Trees) framework. Rather than using a separate validation set, CCP employs a
penalized error criterion that balances training accuracy against tree size, and it generates
an entire sequence of progressively simpler trees from which the best can be selected.

The key idea is to define, for any subtree $T$ of the fully grown tree $T_{\max}$, a
\emph{cost-complexity measure}:
\begin{myequation}
R_\alpha(T) = R(T) + \alpha \lvert \tilde{T} \rvert,
\end{myequation}
where $R(T)$ is the resubstitution error of $T$ (the total misclassification cost on the
training set, possibly weighted by misclassification costs), $\lvert \tilde{T} \rvert$ is
the number of leaf nodes of $T$, and $\alpha \geq 0$ is a complexity parameter that controls
the penalty per leaf. When $\alpha = 0$, no penalty is imposed and $R_0(T)$ is minimized by
the fully grown tree; as $\alpha$ increases, smaller trees become preferred because each
additional leaf must justify its existence by reducing training error by at least $\alpha$.

For a given $\alpha$, the optimal subtree $T(\alpha)$ is the smallest subtree of $T_{\max}$
that minimizes $R_\alpha$. A fundamental result in the CART theory is that the optimal
subtrees form a finite nested sequence:
\begin{myequation}
T_{\max} = T_0 \supset T_1 \supset T_2 \supset \cdots \supset T_m = \{t_1\},
\end{myequation}
where $\{t_1\}$ is the root node alone (the trivial single-leaf tree). This sequence can be
computed efficiently by repeatedly identifying and collapsing the \emph{weakest link}: the
internal node $t$ whose removal produces the smallest increase in resubstitution error per
leaf removed. More precisely, for each internal node $t$ one defines the effective
complexity parameter
\begin{myequation}
\alpha(t) = \frac{R(t) - R(T_t)}{\lvert \tilde{T}_t \rvert - 1},
\end{myequation}
where $R(t)$ is the resubstitution error of the leaf that would replace the subtree $T_t$
rooted at $t$, and $\lvert \tilde{T}_t \rvert$ is the number of leaves in $T_t$. The node
with the smallest $\alpha(t)$ is pruned first, and the process is repeated to generate the
full sequence. The practical implication is that CCP avoids the need to specify $\alpha$ in
advance: the algorithm generates all candidate trees, and cross-validation is then used to
select the best value of $\alpha$ (and hence the best tree in the sequence). This makes CCP
both theoretically principled and practically flexible.

An important feature of CART that distinguishes it from earlier decision tree methods is
that it builds \emph{binary} trees: every internal node has exactly two children, produced
by a binary split on a single feature. This restriction simplifies the tree structure and
the pruning algorithm, and it facilitates the handling of both categorical and continuous
features within a unified framework. For continuous features, binary splits of the form
$x_i \leq \theta$ are used as discussed earlier. For categorical features with $m$ values,
CART considers all $2^{m-1} - 1$ binary partitions of the set of values, which is feasible
for small $m$ but can be expensive for features with high cardinality. CART handles
regression through a separate impurity criterion (the within-node variance, leading to the
squared-error impurity), but the pruning mechanism remains essentially the same.

\subsubsection*{Pessimistic Error Pruning (PEP) and the C4.5 framework}

Pessimistic Error Pruning is the pruning strategy used in the C4.5 algorithm, which itself is an extension of the earlier ID3 algorithm. Unlike REP, PEP does
not require a separate validation set: it estimates the generalization error directly from
the training data using a statistical upper confidence bound on the true error rate. This
makes C4.5 particularly attractive in settings where available data are limited.

The motivation for PEP begins with the observation that the resubstitution error of a
leaf---the fraction of training instances it misclassifies---is an optimistically biased
estimate of the true error. A leaf that correctly classifies all of its training instances
does not necessarily generalize perfectly; it may simply have seen too few instances for its
error rate to be estimated reliably. PEP corrects for this bias by applying a continuity
correction and then using the upper bound of a confidence interval for the binomial
proportion.

Specifically, suppose a leaf contains $n$ instances and misclassifies $e$ of them. Rather
than using $e/n$ as the error estimate, C4.5 computes a pessimistic estimate by treating
the observed error rate as a binomial proportion and computing an upper confidence bound at
a user-specified confidence level (typically $25\%$, a deliberately conservative choice
that encourages pruning). Denoting this upper bound by $\bar{e}(n, e)$, the estimated error
of the leaf is $n \cdot \bar{e}(n, e)$.

The same pessimistic estimate is computed for any subtree rooted at an internal node by
summing the estimated errors of all its leaves. Pruning is then applied top-down: a subtree
is replaced by a leaf if the estimated error of the leaf does not exceed the sum of
estimated errors of the subtree's leaves. Because the upper confidence bound penalizes small
leaf sizes (a leaf with few instances has a wide confidence interval, and hence a higher
upper bound), PEP naturally discourages very fine-grained splits supported by little data
while preserving splits that are well-supported by many instances.

C4.5 also differs from CART in several other respects. It builds multiway splits for
categorical features rather than binary splits, assigns one branch per distinct value of
the feature. It uses information gain ratio as the splitting criterion rather than Gini
impurity, specifically to correct the bias of raw information gain toward high-cardinality
features. It handles missing values through a fractional instance mechanism: when a training
instance has a missing value for the splitting feature, it is distributed fractionally among
all branches in proportion to the fraction of known instances that follow each branch. These
design choices make C4.5 particularly robust in practice and have contributed to its
widespread adoption. The subsequent C5.0 algorithm, a commercial successor to C4.5,
introduced further improvements including boosting integration and more efficient memory
use, but retains the same fundamental pruning philosophy.

\subsubsection*{Minimum Error Pruning (MEP)}

Minimum Error Pruning takes a Bayesian perspective on the
pruning problem. Rather than estimating error through confidence bounds or a validation set,
MEP uses a prior distribution over class probabilities to compute an expected error rate for
each node, and prunes by comparing these expected rates.

The key quantity is the \emph{m-estimate} (or Laplace-corrected estimate) of the error
probability at a node. For a node containing $n$ instances with $n_c$ belonging to the
majority class, the m-estimate of the probability of correct classification is
\begin{myequation}
\hat{p}_c = \frac{n_c + m \cdot p_c}{n + m},
\end{myequation}
where $m$ is a smoothing parameter (often set to the number of classes $k$) and $p_c = 1/k$
is the prior probability of the majority class under a uniform prior. The corresponding
expected error rate is $1 - \hat{p}_c$. For small $n$, the estimate is pulled toward the
prior $1/k$, reflecting the uncertainty due to the small sample; for large $n$, it
approaches the empirical error rate $1 - n_c/n$.

Pruning proceeds bottom-up: a subtree rooted at an internal node is replaced by a leaf if
the expected error of the leaf is no greater than the average expected error of the subtree
(weighted by the number of instances reaching each leaf). In this way, MEP prunes subtrees
that are supported by too few instances to reliably improve upon the prior estimate. A
desirable property of MEP is that it does not require any parameter tuning beyond the
smoothing constant $m$, and it is computationally efficient since it requires only a single
bottom-up pass. Its main limitation is sensitivity to the choice of $m$: different values of
$m$ can lead to substantially different levels of pruning, and there is no universally agreed
upon method for choosing $m$ automatically.

\subsubsection*{Error-Based Pruning (EBP)}

Error-Based Pruning is conceptually related to PEP but differs in that it applies the
confidence-bound correction subtree by subtree rather than leaf by leaf. In EBP, the
estimated error of a subtree is computed by aggregating the confidence-bound-corrected
errors of its individual leaves, and this is compared against the confidence-bound-corrected
error of the corresponding leaf that would result from pruning. The subtree is pruned if the
leaf's estimated error is not substantially worse than the subtree's.

The advantage of EBP over PEP is that it accounts more carefully for the error at the
subtree level, rather than approximating the subtree error as a simple sum of leaf errors.
In practice, EBP tends to produce trees of intermediate size between REP (which can prune
aggressively when the validation set is large) and PEP (which can be conservative when the
confidence level is high). Its main limitation is that, like PEP, it relies on the
confidence level as a hyperparameter whose optimal value is data-dependent.

\subsubsection*{Minimum Description Length (MDL) pruning}

A quite different perspective on pruning comes from information theory, specifically from
the \emph{Minimum Description Length} (MDL) principle. The MDL principle states that the
best model for a dataset is the one that minimizes the total number of bits needed to
describe both the model and the data given the model. In the context of decision trees, this
translates into a trade-off between the description cost of the tree itself (which grows
with the number of nodes) and the description cost of the training labels given the tree
(which decreases as the tree fits the data better).

Concretely, the MDL score of a tree $T$ on dataset $S$ is
\begin{myequation}
MDL(T, S) = L(T) + L(S \mid T),
\end{myequation}
where $L(T)$ is the code length needed to describe the tree structure and its associated
split rules, and $L(S \mid T)$ is the code length needed to describe the training labels
given that the tree is known. Both terms are expressed in bits using appropriate coding
schemes. The pruning algorithm selects the tree that minimizes this combined description
length, balancing model complexity against fitting accuracy. A subtree is pruned if doing
so reduces the total MDL score, meaning that the simplification of the model description
more than compensates for the increased cost of encoding the residual errors.

MDL pruning has attractive theoretical properties: it provides a coherent framework for
model selection grounded in information theory and does not require an arbitrary confidence
parameter. However, it requires careful specification of the coding scheme for the tree and
the labels, and different coding choices can lead to different pruning decisions. In
practice, MDL-based methods are less commonly used in off-the-shelf implementations than
REP or CCP, but they remain an important conceptual tool for understanding the pruning
problem as model selection.

\subsection*{Comparison of pruning methods and practical guidelines}

Having surveyed the main pruning methods, it is instructive to compare them along several
practical dimensions. REP is the most direct method: it optimizes validation accuracy
exactly and is easy to implement, but it requires a dedicated pruning set and may be
unreliable when data are scarce. CCP, as implemented in CART, produces an elegant sequence
of candidate trees and uses cross-validation for model selection, making it robust and
widely applicable; it is the default strategy in the widely used scikit-learn implementation
of decision trees. PEP and EBP, as used in C4.5 and related systems, avoid the need for a
separate validation set at the cost of introducing a confidence parameter, and they are
particularly well-suited to settings with limited data. MEP provides a Bayesian
regularization perspective that is theoretically appealing but requires tuning of the
smoothing parameter. MDL pruning offers a theoretically grounded criterion but requires
additional choices about coding schemes.

Beyond the choice of pruning method, the structural differences between major
implementations also matter in practice. CART builds strictly binary trees with a
unified treatment of regression and classification, uses Gini impurity as the default
splitting criterion, and supports cost-sensitive classification through weighted
misclassification costs. C4.5 builds multiway splits for categorical features, uses
gain ratio to correct for high-cardinality bias, and handles missing values through
fractional instance propagation. ID3, the ancestor of C4.5, uses raw information gain
and does not handle continuous features or missing values natively. More recent
implementations such as LightGBM and XGBoost, while primarily associated with gradient
boosting, incorporate sophisticated variants of decision tree construction and pruning that
include leaf-wise growth strategies, histogram-based split finding, and explicit
regularization terms in the splitting criterion, pushing the performance envelope further
while retaining the interpretability benefits of tree-based models at the individual tree
level.

In most practical applications, the choice of pruning strategy matters less than ensuring
that some form of regularization is applied. An unpruned tree almost always overfits, while
a reasonably pruned tree---regardless of the specific method used---typically achieves
substantially better generalization. When data are plentiful, REP or CCP with
cross-validation are natural first choices. When data are scarce, PEP or MEP avoid the cost
of setting aside a validation set. When theoretical interpretability of the pruning criterion
is important, MDL provides a principled framework. And when the decision tree is used as a
component within an ensemble method such as random forests or gradient boosting, pruning is
often replaced by limiting the maximum depth directly, since the ensemble mechanism provides
its own form of variance reduction that makes aggressive pruning unnecessary at the
individual tree level.


%\end{document}